% ======================================================================
% Appendix E. Consistency checks and verification protocol
% PATH: src/appendices/E-consistency-checks.tex
% ======================================================================

\appendix
\chapter{Consistency checks, invariance audits, and verification protocols}
\label{app:E.consistency}

\section{Purpose and philosophy}
\label{sec:E.purpose}

The localized trace formalism developed in this monograph is designed to be \emph{stable under admissible
modifications} (cutoff changes, quantization choices, microlocal refinements, and normalization conventions) and to
encode all noncanonical dependence in explicit correction terms and a consolidated error ledger.

This appendix records a verification layer: a family of \emph{consistency checks} that can be applied to any concrete
implementation of the localized trace identity.  The goal is not to re-prove the main theorems, but to provide a
portable protocol ensuring that (i) the stated identities are internally consistent, (ii) invariance claims are
properly tracked, and (iii) numerics or symbolic checks are synchronized with the analytic ledger.

The checks are organized as \emph{audits}.  Each audit isolates a structural invariant, states a verifiable criterion,
and lists failure modes.  The protocols are stated in a model-independent way and can be instantiated both in purely
analytic proofs and in computational verification.

\section{Objects and invariants under audit}
\label{sec:E.objects}

We recall the principal objects that enter the trace formalism.

\begin{itemize}
\item The raw localized trace functional $\mathcal{L}_\Gamma[h]$ and its normalized version
\[
\mathcal{T}_\Gamma[h;\chi]=\mathcal{L}_\Gamma[h-h(0)\chi],
\]
with $\chi\equiv 1$ near $0$.

\item The admissible data triple $(\chi,\Op,\Pi)$: cutoff profile $\chi$, quantization scheme $\Op$, and microlocal
partition/projector system $\Pi$.

\item The consolidated ledger
\[
\mathcal{E}_\Gamma[h;\chi]
=
\mathcal{E}^{\mathrm{quant}}_\Gamma[h;\chi]
+
\mathcal{E}^{\mathrm{ref}}_\Gamma[h;\chi]
+
\mathcal{E}^{\mathrm{tail}}_\Gamma[h;\chi]
+
\mathcal{E}^{\mathrm{cont}}_\Gamma[h;\chi],
\]
as introduced in Chapter~\ref{chap:08-normalization}.

\item The trace-to-determinant interface, encoded by oscillatory test families
$h_{\tau,\Delta}(t)=W(t/\Delta)e^{-i\tau t}$ and observables
$\mathfrak{D}(\tau;\Delta)=\mathcal{T}_\Gamma[h_{\tau,\Delta}]$ (Chapter~\ref{chap:10-trace-to-zeta}).

\end{itemize}

The audits below verify invariance properties relative to these objects.

\section{Audit A: transform and normalization conventions}
\label{sec:E.audit.transform}

\subsection{Fourier normalization check}
\label{subsec:E.fourier}

The trace identity depends on the Fourier normalization chosen in Chapters~\ref{chap:02-preliminaries}--\ref{chap:04-wave}.
To avoid hidden sign or factor errors, one must verify the following.

\begin{audit}[A1: Fourier normalization]
\label{audit:E.A1}
\textbf{Invariant.} The pairing between a time test function $h(t)$ and its spectral transform
$\widehat{h}(\xi)$ used in the wave kernel representation is consistent across the entire document.

\textbf{Check.} Verify that:
\begin{enumerate}
\item Plancherel holds in the stated normalization:
\[
\|h\|_{L^2(\R)}^2 = c_\F \|\widehat{h}\|_{L^2(\R)}^2,
\]
with the same constant $c_\F$ as in Chapter~\ref{chap:02-preliminaries}.
\item Evenness is preserved:
if $h(t)=h(-t)$ then $\widehat{h}(\xi)$ is real and even.
\item Convolution squares are nonnegative on the spectral side:
\[
\widehat{h*h^\ast}(\xi)=|\widehat{h}(\xi)|^2\ge 0.
\]
\end{enumerate}

\textbf{Failure modes.} Incorrect $2\pi$ factors; sign mismatch in the phase $e^{-i\xi t}$ versus $e^{+i\xi t}$; loss of
evenness; violation of the quadratic positivity identity.
\end{audit}

\subsection{Reference profile normalization check}
\label{subsec:E.reference}

\begin{audit}[A2: reference profile normalization]
\label{audit:E.A2}
\textbf{Invariant.} The normalization map $\mathsf{N}_\chi(h)=h-h(0)\chi$ removes identity-sector instability and
introduces only deterministic, explicitly computable correction terms.

\textbf{Check.} For two admissible profiles $\chi_1,\chi_2$ with $\chi_i\equiv 1$ near $0$, verify the exact identity
\[
\mathcal{T}_\Gamma[h;\chi_1]-\mathcal{T}_\Gamma[h;\chi_2]
=
-h(0)\,\mathcal{L}_\Gamma[\chi_1-\chi_2].
\]
Confirm that $\mathcal{L}_\Gamma[\chi_1-\chi_2]$ can be evaluated from the identity contribution alone (no hidden
hyperbolic or parabolic dependence).

\textbf{Failure modes.} Dependence on $\chi$ beyond $h(0)$; $\chi$ leaking into hyperbolic classes; nonlocal dependence
on $\chi$ due to an incorrect cutoff hypothesis.
\end{audit}

\section{Audit B: microlocal invariance and quantization stability}
\label{sec:E.audit.microlocal}

\subsection{Quantization change audit}
\label{subsec:E.quantization}

\begin{audit}[B1: change of quantization scheme]
\label{audit:E.B1}
\textbf{Invariant.} Switching between admissible quantizations $\Op$ and $\Op'$ affects the trace functional only
through a controlled ledger term $\mathcal{E}^{\mathrm{quant}}$.

\textbf{Check.} For fixed $(h,\chi,\Pi)$ verify that
\[
\mathcal{T}_\Gamma[h;\chi]_{\Op}
-
\mathcal{T}_\Gamma[h;\chi]_{\Op'}
=
\Err^{\mathrm{quant}}_\Gamma(h;\chi),
\qquad
|\Err^{\mathrm{quant}}_\Gamma(h;\chi)|
\le
\mathcal{E}^{\mathrm{quant}}_\Gamma[h;\chi].
\]
Confirm that the bound depends only on finitely many symbol seminorms and on the localization scale.

\textbf{Failure modes.} Appearance of new leading terms; loss of power-saving; dependence on nonlocal symbol features;
failure of symbolic calculus assumptions (e.g.\ nonuniformity in the semiclassical parameter).
\end{audit}

\subsection{Refinement audit}
\label{subsec:E.refinement}

\begin{audit}[B2: microlocal refinement invariance]
\label{audit:E.B2}
\textbf{Invariant.} Refining the microlocal cover or partition changes the trace only by a controlled refinement
ledger term.

\textbf{Check.} For two admissible microlocal systems $\Pi,\Pi'$ verify
\[
\mathcal{T}_\Gamma[h;\chi]_{\Pi}
-
\mathcal{T}_\Gamma[h;\chi]_{\Pi'}
=
\Err^{\mathrm{ref}}_\Gamma(h;\chi),
\qquad
|\Err^{\mathrm{ref}}_\Gamma(h;\chi)|
\le
\mathcal{E}^{\mathrm{ref}}_\Gamma[h;\chi].
\]
Verify that refinement errors decrease with increased resolution (in the regime where the parametrix is valid).

\textbf{Failure modes.} Nonlocal leakage between microlocal charts; loss of cancellation due to incompatible cutoffs;
accidental duplication/omission of phase space regions.
\end{audit}

\section{Audit C: spectral side integrity (discrete/continuous separation)}
\label{sec:E.audit.spectral}

\subsection{Discrete/continuous decomposition check}
\label{subsec:E.disccont}

In noncompact settings, failure to correctly separate discrete and continuous spectrum is a primary source of hidden
errors.

\begin{audit}[C1: discrete/continuous split]
\label{audit:E.C1}
\textbf{Invariant.} The spectral side decomposes as
\[
\Tr P = \Tr_{\mathrm{disc}} P + \Tr_{\mathrm{cont}} P,
\]
with $\Tr_{\mathrm{cont}} P$ expressed in terms of scattering data, and all dependence on cusp geometry isolated into
explicit terms or the continuous ledger $\mathcal{E}^{\mathrm{cont}}$.

\textbf{Check.} Verify that:
\begin{enumerate}
\item The continuous contribution is computed using the same Plancherel measure and scattering normalization as in
Chapter~\ref{chap:06-parabolic}.
\item The subtraction/normalization scheme does not introduce spurious poles in the continuous integral.
\item The resulting remainder bound satisfies
\[
|\Tr_{\mathrm{cont}} P|\le C\,\mathcal{E}^{\mathrm{cont}}_\Gamma[h;\chi]
\]
in the intended scaling regime.
\end{enumerate}

\textbf{Failure modes.} Missing Maass--Selberg correction; incorrect measure factors; failure of analytic continuation;
continuous terms contaminating the identity main term.
\end{audit}

\section{Audit D: identity term and effective volume synchronization}
\label{sec:E.audit.identity}

The identity sector is the anchor for all later asymptotics.  In localized settings, the identity term is not merely
$\vol(M)$ times a global density; it is an \emph{effective volume} determined by the localization scale and the chosen
normalizations.

\begin{audit}[D1: identity term coherence]
\label{audit:E.D1}
\textbf{Invariant.} The identity contribution computed in Appendix~\ref{app:A.identity} matches the main term used in
all applications (local Weyl laws, variance identities, QUE bounds) under the same normalization conventions.

\textbf{Check.} Verify that:
\begin{enumerate}
\item The leading identity term for the basic projector $P_{\lambda,\eta}$ equals
\[
\frac{\vol(M)}{2\pi}\,\lambda\,\eta
\quad
\text{(in the $\lambda$-parameterization used in Chapter~\ref{ch:applications}).}
\]
\item The same coefficient appears after any admissible cutoff change, with only deterministic correction terms.
\item Any effective-volume modification due to cusp truncation or microlocal exclusion is either explicit or charged
to the ledger.
\end{enumerate}

\textbf{Failure modes.} Inconsistent parameterization ($r$ vs.\ $\lambda$); missing Jacobian factors; hidden cutoff
dependence; mismatch between local and global density constants.
\end{audit}

\section{Audit E: remainder hierarchy and ledger closure}
\label{sec:E.audit.ledger}

\subsection{Ledger decomposition audit}
\label{subsec:E.ledger-decomp}

\begin{audit}[E1: ledger completeness]
\label{audit:E.E1}
\textbf{Invariant.} All noncanonical dependences and secondary analytic terms are absorbed into the consolidated ledger
$\mathcal{E}_\Gamma[h;\chi]$.

\textbf{Check.} For each application (local Weyl, Kuznetsov variance, QE, $L$-moments), enumerate all error sources and
verify that each is assigned to exactly one component:
\[
\mathcal{E}^{\mathrm{quant}},
\quad
\mathcal{E}^{\mathrm{ref}},
\quad
\mathcal{E}^{\mathrm{tail}},
\quad
\mathcal{E}^{\mathrm{cont}}.
\]
Confirm that no term is counted twice and no term remains uncharged.

\textbf{Failure modes.} Double counting between tail and refinement; continuous contributions hidden in ``tail''; loss
of power-saving due to an untracked logarithmic term.
\end{audit}

\subsection{Power-saving closure audit}
\label{subsec:E.powersaving}

\begin{audit}[E2: power-saving closure]
\label{audit:E.E2}
\textbf{Invariant.} The application-scale remainders preserve a strict hierarchy: main term $\gg$ remainder in the
intended regime (e.g.\ $\eta\ge \lambda^{-\theta}$).

\textbf{Check.} Verify that each theorem is stated with a remainder bound of the form
\[
O(\lambda^{1-\delta})
\quad\text{or}\quad
O(\lambda^{-\delta})
\]
consistent with the scaling of the main term.
When an application introduces additional averaging (e.g.\ second moments), verify that the normalization by the local
Weyl count does not lose the power-saving.

\textbf{Failure modes.} Remainder comparable to the main term in short windows; hidden dependence on $\eta$ that breaks
uniformity; loss of $\delta$ in arithmetic weights.
\end{audit}

\section{Audit F: Tauberian interface integrity}
\label{sec:E.audit.tauberian}

The local Weyl and moment statements often require removing smoothing (Appendix~\ref{app:D.tauberian}).
This step must be audited independently.

\begin{audit}[F1: smoothing removal]
\label{audit:E.F1}
\textbf{Invariant.} The Tauberian removal step introduces only a boundary-layer error controlled by monotonicity and
does not affect the trace formula ledger.

\textbf{Check.} Verify that the sharp recovery uses a bracketing inequality of the form
\[
N(r-\delta)\le (N*w_\delta)(r)\le N(r+\delta),
\]
and that the increment term $N(r+\delta)-N(r-\delta)$ is of lower order under the chosen scale $\delta=\delta(r)$.

\textbf{Failure modes.} Use of nonpositive kernels; hidden dependence of $\delta$ on spectral data; circularity (using
the sharp Weyl law to prove itself).
\end{audit}

\section{Audit G: positivity and quadratic-form checks}
\label{sec:E.audit.positivity}

Positivity mechanisms (Li--Weil type criteria) are structurally sensitive to normalization and to the use of
convolution squares.

\begin{audit}[G1: quadratic-form positivity]
\label{audit:E.G1}
\textbf{Invariant.} For $h^\ast(t)=\overline{h(-t)}$, the quadratic form
\[
\mathcal{Q}[h]:=\mathcal{T}_\Gamma[h*h^\ast]
\]
is nonnegative on the spectral side and has its geometric side controlled by the ledger and explicit corrections.

\textbf{Check.} Verify that:
\begin{enumerate}
\item $\widehat{h*h^\ast}(\xi)=|\widehat{h}(\xi)|^2\ge 0$ under the fixed Fourier convention.
\item The normalization subtraction does not break positivity (identity correction is deterministic and accounted).
\item Any truncation/tail in the convolution is charged to $\mathcal{E}^{\mathrm{tail}}$.
\end{enumerate}

\textbf{Failure modes.} Sign errors in Fourier transform; non-even kernels; normalization applied inconsistently inside
the convolution square.
\end{audit}

\section{Audit H: implementation checklist for computational verification}
\label{sec:E.implementation}

When numerical experiments are used (e.g.\ to validate remainder sizes or to sanity-check oscillatory transforms),
one must enforce a strict synchronization between implementation parameters and analytic assumptions.

\subsection{Parameter ledger}
\label{subsec:E.param-ledger}

\begin{audit}[H1: parameter ledger]
\label{audit:E.H1}
\textbf{Invariant.} Every numerical run is fully determined by a finite parameter set that can be recorded and
reproduced.

\textbf{Checklist.} Record:
\begin{itemize}
\item geometry input (group $\Gamma$, cusp widths, truncation height if any);
\item spectral window parameters $(\lambda,\eta)$ and the admissible range checks;
\item cutoff profile $\chi$ and the list of seminorms used in error bounds;
\item Fourier convention and discretization grid (time/spectral);
\item microlocal partition and quantization scheme (if implemented numerically);
\item tail truncation thresholds (geodesic length cut, time cut, frequency cut).
\end{itemize}

\textbf{Failure modes.} Unrecorded implicit parameters; mixed Fourier conventions across modules; tail truncation not
charged to the ledger.
\end{audit}

\subsection{Invariance stress tests}
\label{subsec:E.stress}

\begin{audit}[H2: invariance stress tests]
\label{audit:E.H2}
\textbf{Invariant.} Admissible modifications produce outputs consistent with the deterministic corrections and with
the ledger bounds.

\textbf{Tests.} Perform the following perturbations and verify stability:
\begin{enumerate}
\item Replace $\chi$ by $\chi'$ with the same normalization near $0$ and check that the difference is exactly
$-h(0)\mathcal{L}_\Gamma[\chi-\chi']$ plus a small numerical residual consistent with discretization.
\item Switch quantization (Weyl $\leftrightarrow$ Kohn--Nirenberg) in a controlled symbolic model and verify the
difference is within the predicted $\mathcal{E}^{\mathrm{quant}}$.
\item Refine the microlocal cover and verify convergence within $\mathcal{E}^{\mathrm{ref}}$.
\item Increase tail cutoffs and confirm monotone stabilization within $\mathcal{E}^{\mathrm{tail}}$.
\end{enumerate}

\textbf{Failure modes.} Drift larger than ledger prediction; sign flips; lack of convergence under refinement.
\end{audit}

\section{Summary: the verification layer as a portability guarantee}
\label{sec:E.summary}

The localized trace formalism is intended to be a \emph{portable analytic instrument}.  Portability requires a
verification layer that can detect inconsistencies before they propagate into applications.
The audits in this appendix provide such a layer: each invariance claim is paired with a concrete criterion, and each
criterion can be checked either analytically or computationally.

When these audits pass, the localized trace functional $\mathcal{T}_\Gamma[h;\chi]$ can be treated as a stable object,
and the error ledger $\mathcal{E}_\Gamma[h;\chi]$ becomes a reliable quantitative interface between microlocal trace
analysis and applications in spectral geometry, quantum chaos, and automorphic $L$-function theory.

\section*{Bibliographic notes}

Consistency and invariance checks of the above type are implicit in classical presentations of trace formulae and
microlocal analysis, but they are rarely isolated as explicit audits.
The present appendix adopts an ``audit'' format to make the dependence on conventions and auxiliary choices fully
transparent and to facilitate reproducible implementations.

% ======================================================================
% End of E-consistency-checks.tex
% ======================================================================
